%% ma: L12-B19-0040
%% lop: 12
%% chuong: 6
%% bai: 19
%% dang: 12.19.17
%% loai: DS
%% muc: [NB, TH, TH, VD]
%% dap_an: "ĐSĐĐ"
%% nguon: "Ảnh giáo viên gửi 24/10/2026 (M06): Kỳ thi thử đợt IV, Thi thử TN THPT 2027 (Biên soạn: Aimbet), Phần 2 Câu 3"
%% kien_thuc: "Bài 18 (xác suất có điều kiện), Bài 19 (công thức xác suất toàn phần, công thức Bayes)"
%% ghi_chu: "Đáp án tự giải, đề không kèm đáp án"
%% trang_thai: chinh-thuc
\begin{ex}
Trong một khu rừng, người ta khảo sát hai loài nấm A và B có họ hàng rất gần nhau. Thống kê trên thực tế cho thấy có $70\%$ số nấm trong khu rừng thuộc loài nấm A. Sau khi quan sát có tồn tại hay không vòng nằm trên cuống của cây nấm, người ta thống kê rằng tỉ lệ cây nấm có đặc điểm này ở loài nấm A và loài nấm B lần lượt là $80\%$ và $30\%$. Chọn ngẫu nhiên một cây nấm tại khu rừng.
\choiceTF
{\True Biết rằng cây nấm là loài nấm A, xác suất để cây nấm \textbf{không} có vòng nằm trên cuống bằng $20\%$.}
{Biết rằng cây nấm có vòng nằm trên cuống, xác suất để cây nấm là loài nấm B lớn hơn $15\%$.}
{\True Xác suất để cây nấm có vòng nằm trên cuống bằng $65\%$.}
{\True Nếu người ta phân loại rằng các cây nấm có vòng là loài nấm A và các cây nấm \textbf{không} có vòng là loài nấm B, xác suất để cây nấm bị phân loại sai so với thực tế nhỏ hơn $25\%$.}
\loigiai{Gọi $A$, $B$: cây nấm thuộc loài A, B; $V$: có vòng. $P(A)=0{,}7$, $P(B)=0{,}3$, $P(V\mid A)=0{,}8$, $P(V\mid B)=0{,}3$.
a) $P\left(\overline V\mid A\right)=0{,}2$: đúng.
c) $P(V)=0{,}7\cdot0{,}8+0{,}3\cdot0{,}3=0{,}65$: đúng.
b) $P(B\mid V)=\dfrac{0{,}09}{0{,}65}\approx0{,}138<0{,}15$: sai.
d) Phân loại sai khi nấm A không có vòng hoặc nấm B có vòng: $0{,}7\cdot0{,}2+0{,}3\cdot0{,}3=0{,}23<0{,}25$: đúng.}
\end{ex}
